\chapter{Компактно-фазовый гейт embedding-нормировки}

Предыдущий гейт показал, что числовая несовместимость двух норм отсутствует,
но единичная амплитуда лептонной нуль-формы не инвариантна относительно
непрерывной замены координаты. Минимальная попытка устранить этот разрыв
состоит в замене вещественной амплитуды на компактную фазу
\begin{equation}
 U=e^{i\vartheta},
 \qquad
 \vartheta\sim\vartheta+2\pi.
\end{equation}

\section{Что компактность действительно фиксирует}

Период \(2\pi\) запрещает произвольную непрерывную замену
\(\vartheta\mapsto a\vartheta\): она не является автоморфизмом \(U(1)\),
кроме \(a=\pm1\). Кроме того, локальные вершины имеют вид
\(U^q=e^{iq\vartheta}\) с целым зарядом \(q\). Поэтому при заранее
фиксированном примитивном заряде \(q=1\) линейная нормировка фазовой
координаты больше не является свободным непрерывным параметром.

\begin{proposition}[Кинематический embedding-pass]
Если лептонная коллективная координата является примитивной фазой
компактной группы \(U(1)\), её масштаб фиксирован периодом и решёткой
зарядов с точностью до знака. Разрыв непрерывного field redefinition
\(a_\tau\) из предыдущей главы на кинематическом уровне исчезает.
\end{proposition}

Это утверждение фиксирует координату, но ещё не фиксирует метрику на
пространстве фаз и не гарантирует ненулевую жёсткость постоянной моды.

\section{Минимальное локальное действие}

Для sigma-model действия на связном носителе \(Y\)
\begin{equation}
 S_{\mathrm{kin}}[\vartheta]
 =\frac{f^2}{2}\int_Y d\vartheta\wedge\star d\vartheta
\end{equation}
постоянная вариация лежит в ядре:
\begin{equation}
 \Hess_{\vartheta_0}S_{\mathrm{kin}}[1,1]=0.
\end{equation}
Следовательно, объёмный вклад \(\Vol(Y)\) не возникает из
производного действия. Коэффициент \(f^2\) также является радиусом
групповой метрики, а не следствием одного лишь периода фазы.

Ненулевой гессиан можно получить из периодического потенциала
\begin{equation}
 S_{\mathrm{pot}}[\vartheta]
 =\mu^{\dim Y}\int_Y(1-\cos\vartheta)\,d\Vol_Y,
\end{equation}
но тогда
\begin{equation}
 \Hess_0S_{\mathrm{pot}}[1,1]
 =\mu^{\dim Y}\Vol(Y).
\end{equation}
Компактность фиксирует форму косинуса и целочисленный заряд, но не фиксирует
масштаб \(\mu\). Назначение \(\mu=1\) без дополнительного принципа является
новой нормировкой действия.

\section{Локальность на полном носителе}

Одна постоянная фаза на
\(K=\mathbb{RP}^3\times S^1\) даёт объём полного произведения
\begin{equation}
 \Vol(K)=2\pi^3,
\end{equation}
а не сумму \(\pi^2+2\pi\). Требуемая сумма возникает только для
стратифицированного конфигурационного пространства
\begin{equation}
 \operatorname{Map}(\mathbb{RP}^3,U(1))
 \oplus\operatorname{Map}(S^1,U(1)),
\end{equation}
с отдельными интегралами по двум факторам. Такая конструкция алгебраически
даёт
\begin{equation}
 \|1_{\mathbb{RP}^3}\|^2+\|1_{S^1}\|^2=\pi^2+2\pi,
\end{equation}
но её нельзя считать локальным скалярным действием на \(K\), пока две меры
и две embedding-карты не выведены как граничные, дефектные или условные
следы одного parent action.

\begin{nogocriterion}[Минимальный compact-phase no-go]
Производное действие одной компактной фазы на \(K\) имеет нулевой гессиан
на постоянной моде. Периодический потенциал вводит нефиксированный масштаб,
а локальный объём полного носителя даёт произведение, а не требуемую сумму
объёмов факторов. Поэтому минимальный локальный compact-phase вариант не
закрывает двухсекторный parent-action gate.
\end{nogocriterion}

\section{Оставшаяся открытая ветвь}

Компактная фаза не отвергается полностью: она устраняет непрерывную
неоднозначность самой координаты. Открытым остаётся стратифицированный
вариант, в котором две фазовые моды возникают из единого граничного или
дефектного комплекса, а общий коэффициент их гессиана фиксируется той же
канонической pairing form, что и нейтринная норма. Следующим гейтом должна
стать именно конструкция такой pairing form; добавление произвольного
decay constant или двух независимых потенциалов запрещено.

Машинный аудит сохранён в
\texttt{s2t\_v3\_compact\_phase\_embedding\_results.json}.